\section{CoM高さ変動とZMP挙動の関係}

従来の多くの歩行制御手法では，線形倒立振子モデル（LIPM）に基づき，CoMの高さを一定と仮定している．この仮定の下では，CoMの水平方向運動とZMPの関係が単純化され，安定性解析や制御設計が容易となる．しかし，軟弱地面において床沈み込みが生じる場合，この仮定は成立しない．

支持脚側の床が沈み込むと，ロボット全体のCoM高さは低下する．このCoM高さの変動は，ZMPの位置に直接的な影響を与える．一般に，CoM高さが低下すると，同一の水平方向加速度に対してZMPの変位量が変化し，結果としてZMPが支持多角形の境界付近，あるいは外側へと移動しやすくなる．

特に，片脚支持期では支持多角形が小さいため，わずかなZMP変動であっても安定性に大きな影響を及ぼす．そのため，CoM高さ変動を伴う状況では，従来の高さ一定モデルに基づく制御では，十分な安定性を確保できない可能性がある．このことは，軟弱地面における歩行において，CoMの鉛直方向運動を無視できないことを示唆している．

\( ※ 改変の余地あり \)